% GENERATED FILE. Edit canonical lessons, then run npm run generate:books.
% canonical-source-hash: fnv1a64:4cd9fc19125f76c5
% canonical-lessons: ZH-C324-xiangqi, ZH-C324-jinian, ZH-C324-liuxia, ZH-C324-shidai, ZH-C324-yinxiang

\chapter{Remembering}
\label{ch:zh-remembering}
\hlchaptermodality{324}

\begin{chapteropening}
\textbf{By the end of this chapter:} \emph{I can say something came back to me, honour someone's memory, leave something behind and talk about an era.}

Complete the last lesson of chapter 324: I can say something came back to me, honour someone's memory, leave something behind and talk about an era.
\end{chapteropening}

\section[\texorpdfstring{\zh{想起} (xiǎngqǐ)}{想起 (xiǎngqǐ)}]{\zh{想起} (xiǎngqǐ) — to think of; to call to mind}

\label{lesson:ZH-C324-xiangqi}



\begin{quote}
Before the new one: say the Mandarin for ancient times, then the Mandarin for modern times.
\end{quote}

\subsection*{You'll want to know: \zh{想起}}
\textbf{\zh{想起}} — \emph{xiǎngqǐ} — \textquotedblleft{}to think of; to call to mind\textquotedblright{}.

\begin{quote}
\textbf{\zh{我突然想起了他的信}} \emph{wǒ tūrán xiǎngqǐ le tā de xìn} — I suddenly remembered his letter.
\end{quote}

Two third tones meet in this word: xiǎng is said with a rising tone before qǐ, and the spelling keeps both marks.

\begin{cousinweb}
\textbf{\zh{想起}} (\emph{xiǎngqǐ}) is built from \zh{想} 'to think' and \zh{起} 'to rise', both inherited from Middle Chinese, \emph{sjangX} and \emph{khiX} (Guangyun, in Baxter's transcription), which glosses the first 'to think' and the second 'to rise'.
\end{cousinweb}

\begin{grammarlens}[title={What you've built}]
One more word for telling what happened.
\end{grammarlens}

\subsection*{Your turn}
\begin{itemize}
\raggedright
  \item \emph{Say it:} \emph{xiǎngqǐ}
  \item \emph{Say it:} \emph{xiǎngqǐ}, once more
  \item \emph{Say it:} \emph{wǒ tūrán xiǎngqǐ le tā de xìn} — I suddenly remembered his letter.
  \item \emph{Recall:} say the Mandarin for ancient times, then the Mandarin for modern times, then say \emph{xiǎngqǐ} again
\end{itemize}

\begin{tcolorbox}[breakable,colback=teal!4,colframe=teal!35!black,title={Before you move on}]
What does \zh{想起} mean? (\textquotedblleft{}To think of; to call to mind\textquotedblright{}.) Say it once more.
\end{tcolorbox}
\section[jìniàn]{jìniàn — to commemorate; to honour the memory of}

\label{lesson:ZH-C324-jinian}



\begin{quote}
Before the new one: say the Mandarin for modern times, then the Mandarin for to call to mind.
\end{quote}

\subsection*{You'll want to know: jìniàn}
\emph{jìniàn} — \textquotedblleft{}to commemorate; to honour the memory of\textquotedblright{}.

\begin{quote}
\emph{měinián zhè yì tiān wǒmen dōu jìniàn tā} — Every year on this day we remember him.
\end{quote}

Its characters come later; for now it is a word for the ear.

\begin{cousinweb}
\emph{jìniàn} is built from \emph{jì} and \emph{niàn}, both inherited from Middle Chinese, \emph{kiX} and \emph{nemH} (Guangyun, in Baxter's transcription), which glosses the first, among other senses, 'to note' and 'a span of twelve years', and the second 'to think of'.
\end{cousinweb}

\begin{grammarlens}[title={What you've built}]
One more word for telling what happened.
\end{grammarlens}

\subsection*{Your turn}
\begin{itemize}
\raggedright
  \item \emph{Say it:} \emph{jìniàn}
  \item \emph{Say it:} \emph{jìniàn}, once more
  \item \emph{Say it:} \emph{měinián zhè yì tiān wǒmen dōu jìniàn tā} — Every year on this day we remember him.
  \item \emph{Recall:} say the Mandarin for modern times, then the Mandarin for to call to mind, then say \emph{jìniàn} again
\end{itemize}

\begin{tcolorbox}[breakable,colback=teal!4,colframe=teal!35!black,title={Before you move on}]
What does jìniàn mean? (\textquotedblleft{}To commemorate; to honour the memory of\textquotedblright{}.) Say it once more.
\end{tcolorbox}
\section[liúxià]{liúxià — to leave behind}

\label{lesson:ZH-C324-liuxia}



\begin{quote}
Before the new one: say the Mandarin for to call to mind, then the Mandarin for to honour the memory of.
\end{quote}

\subsection*{You'll want to know: liúxià}
\emph{liúxià} — \textquotedblleft{}to leave behind\textquotedblright{}.

\begin{quote}
\emph{yéye gěi wǒmen liúxià le yì běn gǔshū} — Grandpa left us an old book.
\end{quote}

Its characters come later; for now it is a word for the ear.

\begin{cousinweb}
\emph{liúxià} is built from \emph{liú} 'to stay' and \zh{下} 'down', both inherited from Middle Chinese, \emph{ljuw} and \emph{haeX} (Guangyun, in Baxter's transcription), which glosses the first 'to stay' and 'to stop' and the second 'below'.
\end{cousinweb}

\begin{grammarlens}[title={What you've built}]
One more word for telling what happened.
\end{grammarlens}

\subsection*{Your turn}
\begin{itemize}
\raggedright
  \item \emph{Say it:} \emph{liúxià}
  \item \emph{Say it:} \emph{liúxià}, once more
  \item \emph{Say it:} \emph{yéye gěi wǒmen liúxià le yì běn gǔshū} — Grandpa left us an old book.
  \item \emph{Recall:} say the Mandarin for to call to mind, then the Mandarin for to honour the memory of, then say \emph{liúxià} again
\end{itemize}

\begin{tcolorbox}[breakable,colback=teal!4,colframe=teal!35!black,title={Before you move on}]
What does liúxià mean? (\textquotedblleft{}To leave behind\textquotedblright{}.) Say it once more.
\end{tcolorbox}
\section[\texorpdfstring{\zh{时代} (shídài)}{时代 (shídài)}]{\zh{时代} (shídài) — era; age}

\label{lesson:ZH-C324-shidai}



\begin{quote}
Before the new one: say the Mandarin for to honour the memory of, then the Mandarin for to leave behind for others.
\end{quote}

\subsection*{You'll want to know: \zh{时代}}
\textbf{\zh{时代}} — \emph{shídài} — \textquotedblleft{}era; age\textquotedblright{}.

\begin{quote}
\textbf{\zh{现在是手机的时代}} \emph{xiànzài shì shǒujī de shídài} — This is the age of the mobile phone.
\end{quote}

\begin{cousinweb}
\textbf{\zh{时代}} (\emph{shídài}) is built from \zh{时} 'a time' and \zh{代} 'an age', both inherited from Middle Chinese, \emph{dzyi} and \emph{dojH} (Guangyun, in Baxter's transcription), which glosses the first with another word for a time and the second 'to take turns'.
\end{cousinweb}

\begin{grammarlens}[title={What you've built}]
One more word for telling what happened.
\end{grammarlens}

\subsection*{Your turn}
\begin{itemize}
\raggedright
  \item \emph{Say it:} \emph{shídài}
  \item \emph{Say it:} \emph{shídài}, once more
  \item \emph{Say it:} \emph{xiànzài shì shǒujī de shídài} — This is the age of the mobile phone.
  \item \emph{Recall:} say the Mandarin for to honour the memory of, then the Mandarin for to leave behind for others, then say \emph{shídài} again
\end{itemize}

\begin{tcolorbox}[breakable,colback=teal!4,colframe=teal!35!black,title={Before you move on}]
What does \zh{时代} mean? (\textquotedblleft{}Era; age\textquotedblright{}.) Say it once more.
\end{tcolorbox}
\section[yìnxiàng]{yìnxiàng — impression}

\label{lesson:ZH-C324-yinxiang}



\begin{quote}
Before the new one: say the Mandarin for to leave behind for others, then the Mandarin for an age in history.
\end{quote}

\subsection*{You'll want to know: yìnxiàng}
\emph{yìnxiàng} — \textquotedblleft{}impression\textquotedblright{}.

\begin{quote}
\emph{wǒ duì nàge dìfang de yìnxiàng hěn hǎo} — That place made a very good impression on me.
\end{quote}

Its characters come later; for now it is a word for the ear.

\begin{cousinweb}
\emph{yìnxiàng} is built from \emph{yìn} 'a seal' and \emph{xiàng} 'an image', both inherited from Middle Chinese, \emph{'jinH} and \emph{zjangX} (Guangyun, in Baxter's transcription), which glosses the first 'a seal' and describes the second as the long-nosed beast of the south. The sense 'an image' is not given there.
\end{cousinweb}

\begin{grammarlens}[title={What you've built}]
One more word for telling what happened.
\end{grammarlens}

\subsection*{Your turn}
\begin{itemize}
\raggedright
  \item \emph{Say it:} \emph{yìnxiàng}
  \item \emph{Say it:} \emph{yìnxiàng}, once more
  \item \emph{Say it:} \emph{wǒ duì nàge dìfang de yìnxiàng hěn hǎo} — That place made a very good impression on me.
  \item \emph{Recall:} say the Mandarin for to leave behind for others, then the Mandarin for an age in history, then say \emph{yìnxiàng} again
\end{itemize}

\begin{tcolorbox}[breakable,colback=teal!4,colframe=teal!35!black,title={Before you move on}]
What does yìnxiàng mean? (\textquotedblleft{}Impression\textquotedblright{}.) Say it once more.
\end{tcolorbox}
